\chapter{Trunking, 802.1Q and SVIs}\label{ch:trunking}
\bp{2.1}

\begin{outcomes}
By the end of this chapter you will be able to:
\begin{tight}
  \item \textbf{Explain} what an 802.1Q tag is and why the native VLAN is the
        exception that causes most trunk faults.
  \item \textbf{Configure} a trunk between two switches and between a switch and
        a router, and prune it to the VLANs it should carry.
  \item \textbf{Configure} switch virtual interfaces so that a multilayer switch
        routes between VLANs.
  \item \textbf{Diagnose} the four trunk faults: mode mismatch, encapsulation
        mismatch, native VLAN mismatch and a missing VLAN in the allowed list.
\end{tight}
\end{outcomes}

\begin{prereq}
\chref{vlans} for VLANs and access ports; \chref{ipv4} for addressing, which the
SVI section needs.
\end{prereq}

\begin{keyterms}
trunk $\bullet$ 802.1Q $\bullet$ tag $\bullet$ native VLAN $\bullet$ allowed VLAN
list $\bullet$ DTP $\bullet$ dynamic auto $\bullet$ dynamic desirable $\bullet$
SVI $\bullet$ inter-VLAN routing $\bullet$ router on a stick $\bullet$
subinterface
\end{keyterms}

\section{The tag, and the one VLAN that has none}

A single link between two switches must carry traffic for every VLAN, and the
receiving switch has to know which VLAN each frame belongs to. \term{802.1Q}
solves this by inserting a four-byte tag into the frame carrying the VLAN ID.

\ccnafig{%
\begin{tikzpicture}[font=\scriptsize]
  \tikzset{
    dv/.style={rounded corners=3pt, draw=#1, fill=#1!8, text=#1!45!black,
               line width=0.9pt, minimum width=1.45cm, minimum height=1.0cm,
               align=center, font=\tiny\bfseries},
    w/.style={line width=0.85pt, neutral!65},
    fld/.style={draw=#1, fill=#1!12, text=#1!40!black, minimum height=0.6cm,
                align=center, font=\tiny\bfseries, line width=0.7pt}}

  \node[dv=dom2] (s1) at (0,0)   {\faIcon{sitemap}\\SW1};
  \node[dv=dom2] (s2) at (12.6,0) {\faIcon{sitemap}\\SW2};
  \draw[w, line width=1.5pt, dom2!60] (s1) -- (s2);
  \node[font=\tiny\bfseries, text=dom2!45!black] at (6.3,0.42)
      {one physical link --- the \emph{trunk}};

  \foreach \y/\v/\c in {-1.15/10/dom1, -1.75/20/dom4, -2.35/30/dom5} {
    \draw[-{Stealth[length=2mm]}, \c!70, line width=0.9pt] (0.9,\y) -- (11.7,\y);
    \node[font=\tiny\bfseries, text=\c!45!black, anchor=east] at (0.85,\y) {VLAN \v};
    \node[font=\tiny, text=\c!45!black, fill=white, inner sep=1.5pt] at (6.3,\y)
        {tagged \textbf{\v}};}
  \draw[-{Stealth[length=2mm]}, neutral!70, line width=0.9pt] (0.9,-2.95) -- (11.7,-2.95);
  \node[font=\tiny\bfseries, text=neutral!55!black, anchor=east] at (0.85,-2.95) {VLAN 99};
  \node[font=\tiny, text=accent!70!black, fill=white, inner sep=1.5pt] at (6.3,-2.95)
      {\textbf{untagged} --- the native VLAN};

  \node[font=\tiny, text=neutral!85!black, align=center] at (6.3,-3.75)
      {Three VLANs carry a tag that says which they are.\\
       The native VLAN carries none --- so both ends must agree what it is.};
\end{tikzpicture}}%
{A trunk is one cable carrying many VLANs. Every VLAN but one is labelled; the
unlabelled one is the native VLAN, and a disagreement about it joins two
networks silently.}{trunktag}

\begin{conceptbox}[The native VLAN is untagged, and that is the whole problem]
One VLAN on every trunk is designated the \term{native VLAN}, and its frames
cross the trunk \textbf{untagged}. The receiving switch assumes any untagged
frame it sees on a trunk belongs to \emph{its own} native VLAN.

If the two ends disagree about which VLAN is native, traffic silently leaks
between VLANs: frames sent untagged in VLAN 1 by one switch are received into
VLAN 99 by the other. Nothing goes down, nothing logs an error at layer 2, and
two VLANs that should be isolated are now joined. CDP will complain, which is one
of several reasons to leave it on (\chref{discovery}).

Set the native VLAN explicitly and identically at both ends, and make it an
unused VLAN --- never VLAN 1, and never a VLAN carrying users.
\end{conceptbox}

\begin{config}{a trunk between two switches, configured deliberately}
! On BOTH ends, identically.
interface GigabitEthernet0/1
 description Uplink to SW2
 ! Some platforms require the encapsulation before the mode; on
 ! switches that only support dot1q the command is rejected as
 ! unavailable, which is not an error.
 switchport trunk encapsulation dot1q
 switchport mode trunk
 ! An unused VLAN as native, matched at the far end.
 switchport trunk native vlan 999
 ! Carry only what this link needs. Pruning limits broadcast
 ! flooding and limits the blast radius of a mistake.
 switchport trunk allowed vlan 10,20,30
 ! Do not negotiate -- we have decided what this port is.
 switchport nonegotiate
 no shutdown
\end{config}

\begin{verify}{SW1}{show interfaces trunk}
Port        Mode         Encapsulation  Status        Native vlan
Gi0/1       on           802.1q         trunking      999

Port        Vlans allowed on trunk
Gi0/1       10,20,30

Port        Vlans allowed and active in management domain
Gi0/1       10,20,30

Port        Vlans in spanning tree forwarding state and not pruned
Gi0/1       10,20,30
\end{verify}

Four blocks, and they narrow as they go. A VLAN can be \emph{allowed} but not
\emph{active} (it does not exist on this switch), or active but not
\emph{forwarding} (spanning tree is blocking it). When a VLAN is missing from the
last block but present in the first, you have found the layer at which it is
being lost.

\begin{alertbox}[\cmd{allowed vlan} replaces; \cmd{allowed vlan add} appends]
Typing \cmd{switchport trunk allowed vlan 40} on a live trunk does not add
VLAN~40 --- it makes 40 the \emph{only} allowed VLAN and instantly drops every
other VLAN on that link. On a production uplink this is an outage.

Use \cmd{switchport trunk allowed vlan add 40}. The equivalents are
\cmd{remove}, \cmd{except} and \cmd{all}.
\end{alertbox}

\section{DTP, and why to turn it off}

\term{Dynamic Trunking Protocol} lets two Cisco switches negotiate whether a link
becomes a trunk. It is convenient and it is a liability.

\begin{longtable}{@{}L{3.6cm}L{3.4cm}L{3.4cm}L{4.2cm}@{}}
\toprule
\rowcolor{primary!15}
\textbf{Local mode} & \textbf{vs \cmd{trunk}} & \textbf{vs \cmd{dynamic desirable}} & \textbf{vs \cmd{dynamic auto} / \cmd{access}} \\
\midrule
\endhead
\cmd{trunk} & trunk & trunk & trunk / \textbf{mismatch} \\
\rowcolor{lightbg}
\cmd{dynamic desirable} & trunk & trunk & trunk / access \\
\cmd{dynamic auto} & trunk & trunk & \textbf{access} / access \\
\rowcolor{lightbg}
\cmd{access} & \textbf{mismatch} & access & access / access \\
\bottomrule
\end{longtable}

The two entries worth memorising are the failures. \cmd{dynamic auto} on both
ends gives you an \textbf{access port}, not a trunk --- each side waits for the
other to ask. And \cmd{trunk} facing \cmd{access} is a genuine mismatch that
passes only the access VLAN.

\begin{conceptbox}[Configure both ends; negotiate nothing]
Set \cmd{switchport mode trunk} and \cmd{switchport nonegotiate} on uplinks, and
\cmd{switchport mode access} on edge ports. Then a port is what you said it is,
and an attacker who plugs a laptop into a wall socket cannot negotiate himself a
trunk into every VLAN.
\end{conceptbox}

\section{Routing between VLANs}

A VLAN is a broadcast domain and traffic cannot leave it without being routed.
There are two ways to provide that, and the exam names the modern one.

\subsection{Switch virtual interfaces --- the modern way}

An \term{SVI} is a virtual layer 3 interface inside a multilayer switch,
representing one VLAN. Give each VLAN an SVI with an address and the switch
routes between them at hardware speed.

\begin{config}{a multilayer switch routing three VLANs}
! Routing must be switched on -- easily forgotten, and without it the
! SVIs come up and nothing is routed between them.
ip routing
!
interface Vlan10
 description Staff gateway
 ip address 192.168.10.1 255.255.255.0
 no shutdown
!
interface Vlan20
 description Guest gateway
 ip address 192.168.20.1 255.255.255.0
 no shutdown
!
interface Vlan30
 description Voice gateway
 ip address 192.168.30.1 255.255.255.0
 no shutdown
\end{config}

\begin{alertbox}[What it takes to bring an SVI up]
An SVI shows \cmd{up/up} only when \emph{all} of the following hold:
\begin{tight}
  \item the VLAN exists in the VLAN database and is not shut down;
  \item at least one access port in that VLAN is up, \textbf{or} the VLAN is
        allowed and forwarding on a trunk;
  \item the SVI itself is not shut down.
\end{tight}
An SVI that stays down with a perfectly good address almost always means the
second condition: no live port in that VLAN. It is not an addressing fault.
\end{alertbox}

\begin{verify}{SW1}{show ip interface brief | include Vlan}
Vlan10                 192.168.10.1    YES manual up                    up
Vlan20                 192.168.20.1    YES manual up                    up
Vlan30                 192.168.30.1    YES manual up                    down
\end{verify}

VLAN 30's SVI is up administratively and down at protocol level --- no active
port in VLAN 30. Plug in a phone, or allow VLAN 30 on a trunk, and it comes up.

\subsection{Router on a stick --- the older way}

Where the switch cannot route, one router interface carries every VLAN over a
trunk, divided into \term{subinterfaces}.

\begin{config}{router on a stick}
interface GigabitEthernet0/0
 description Trunk to SW1
 no ip address
 no shutdown
!
interface GigabitEthernet0/0.10
 description Staff
 encapsulation dot1Q 10
 ip address 192.168.10.1 255.255.255.0
!
interface GigabitEthernet0/0.20
 description Guest
 encapsulation dot1Q 20
 ip address 192.168.20.1 255.255.255.0
!
! The native VLAN's subinterface is marked, not tagged.
interface GigabitEthernet0/0.999
 encapsulation dot1Q 999 native
\end{config}

Every VLAN's traffic crosses the one physical link twice --- in and out --- so
this scales badly. It appears in the exam because it still appears in small
networks.

\begin{pitfall}
\begin{tight}
  \item \textbf{Native VLAN mismatch.} The two ends disagree, VLANs leak, and
        nothing goes down. CDP logs it; look for
        \cmd{\%CDP-4-NATIVE\_VLAN\_MISMATCH}.
  \item \textbf{Bare \cmd{allowed vlan} on a live trunk.} Replaces the list.
        Always \cmd{add}.
  \item \textbf{Forgetting \cmd{ip routing} on a multilayer switch.} The SVIs
        come up and traffic still will not cross between VLANs.
  \item \textbf{Expecting an SVI to be up with no active port in its VLAN.} It
        will not be, and no amount of re-addressing changes that.
  \item \textbf{Leaving VLAN 1 as native.} It carries control traffic and it is
        every attacker's first guess.
\end{tight}
\end{pitfall}

\begin{breakfix}[One VLAN crosses the uplink; the other three do not.]
After a maintenance window, users in VLAN 10 are fine and users in 20, 30 and 40
cannot reach the servers. Both uplink ports show as trunking.

\begin{verify}{SW2}{show interfaces trunk}
Port        Mode         Encapsulation  Status        Native vlan
Gi0/1       on           802.1q         trunking      999

Port        Vlans allowed on trunk
Gi0/1       10
\end{verify}

\textbf{Diagnosis.} The trunk is healthy; the allowed list has been reduced to
VLAN 10 alone. Somebody ran \cmd{switchport trunk allowed vlan 10} intending to
add VLAN 10 to the existing list, and replaced it instead. The link stayed up
throughout, which is why nothing alerted.

\textbf{Fix.} Restore the full list explicitly rather than adding one VLAN at a
time: \cmd{switchport trunk allowed vlan 10,20,30,40}. Verify with
\cmd{show interfaces trunk} on \emph{both} ends --- the allowed list is
configured per side, and the effective set is the intersection, so a matching
mistake at the far end would leave the fault half-fixed and twice as confusing.
\end{breakfix}

\begin{lab}{Trunk two switches and route between VLANs}{Packet Tracer}
\textbf{Topology.} Two switches trunked together; a multilayer switch or router
providing gateways; two PCs per VLAN across both switches.

\begin{tasks}
  \item Create VLANs 10, 20 and 999 on both switches. Put PCs in 10 and 20 on
        each switch.
  \item Before configuring the trunk, try to ping between the two switches'
        VLAN 10 PCs. Record the failure and explain it.
  \item Configure the trunk on both ends with native VLAN 999 and an allowed list
        of 10, 20 and 999. Re-test.
  \item Change the native VLAN on one end only to 998. Record what still works,
        what breaks, and what the log says.
  \item Restore it, then run \cmd{switchport trunk allowed vlan 20} on one end
        and observe the outage you have just caused. Fix it with \cmd{add}.
  \item Enable \cmd{ip routing} and configure SVIs for VLANs 10 and 20. Prove
        inter-VLAN routing works.
  \item \textbf{Extension.} Shut every access port in VLAN 20 and watch the SVI
        go down. Explain why in terms of the three conditions above.
\end{tasks}
\end{lab}

\begin{checkpoint}
\begin{tightnum}
  \item Both ends of a link are set to \cmd{dynamic auto}. What does the link
        become, and why?
  \item An SVI has a correct address and shows \cmd{up/down}. What is the most
        likely cause and how would you confirm it?
  \item Why is a native VLAN mismatch more dangerous than a trunk that fails to
        come up at all?
\end{tightnum}
\end{checkpoint}

\begin{chaptersummary}
\begin{tight}
  \item 802.1Q tags every VLAN except the native one, which crosses untagged.
        Mismatched native VLANs leak traffic between VLANs silently.
  \item Configure both ends explicitly and add \cmd{switchport nonegotiate}.
        \cmd{dynamic auto} on both ends yields an access port, not a trunk.
  \item \cmd{allowed vlan} replaces the list; \cmd{allowed vlan add} appends.
        This distinction causes real outages.
  \item \cmd{show interfaces trunk} narrows from allowed to active to forwarding;
        the block where a VLAN disappears tells you the layer at fault.
  \item An SVI needs \cmd{ip routing}, an existing unshut VLAN, and at least one
        live port in that VLAN.
\end{tight}
\end{chaptersummary}

\begin{reviewq}
  \item \textbf{(MCQ)} Which command adds VLAN 40 to a trunk without disturbing
        the VLANs already carried?
        (a)~\cmd{switchport trunk allowed vlan 40}
        (b)~\cmd{switchport trunk allowed vlan add 40}
        (c)~\cmd{switchport trunk vlan 40}
        (d)~\cmd{switchport access vlan 40}
  \item \textbf{(MCQ)} A trunk's two ends have different native VLANs. What
        happens? (a)~the trunk fails to form (b)~both ends shut down (c)~untagged
        traffic is bridged between two VLANs (d)~all VLANs are pruned
  \item \textbf{(Short)} Explain why an SVI can show \cmd{up/down} when its
        address and mask are correct.
  \item \textbf{(Short)} Give two reasons to set the native VLAN to an unused
        VLAN rather than leaving it as VLAN 1.
  \item \textbf{(Applied)} A colleague reports that after adding a new VLAN,
        three existing VLANs stopped working across an uplink. State what they
        almost certainly typed, what they should have typed, and how you would
        verify the repair on both ends.
\end{reviewq}
